\clearpage
\begin{proofbox}

\textbf{\textcolor{CProof}{方法一：最低门槛——用整数表看懂规律}}
\begin{center}\begin{tikzpicture}[font=\small,>=stealth,line join=round]
\draw[gray!40,rounded corners=3pt] (0,0) rectangle (14.2,8.8);\node at (7.1,8.200000000000001) {看见每次乘底数：横着走同样一步，高度按倍数改变};
\node at (3.1500000000000004,7.3) {$y=2^x$};\begin{scope}[shift={(1.3758620689655174,3.3000000000000003)},x=1.6896551724137934cm,y=0.6cm]\draw[->,gray] (-0.4,0)--(2.5,0) node[right]{$x$};\draw[->,gray] (0,-0.5)--(0,5) node[above]{$y$};\draw[blue!75!black,thick,solid,domain=-0.2:2.2,samples=180] plot(\x,{pow(2,\x)});\fill[orange!85!black] (0,1) circle(2pt);\fill[orange!85!black] (1,2) circle(2pt);\fill[orange!85!black] (2,4) circle(2pt);\node[below left=5pt] at (0,0){$0$};\node[below=5pt] at (1,0){$1$};\node[below=5pt] at (2,0){$2$};\draw[gray,densely dotted] (1,0)--(1,2);\draw[gray,densely dotted] (2,0)--(2,4);\draw[->,blue!75!black,thick] (.2,1.148698354997035)--(.8,1.7411011265922482);\draw[->,blue!75!black,thick] (1.2,2.2973967099940698)--(1.8,3.4822022531844965);\end{scope}\node[align=center,text width=5.8cm] at (3.1500000000000004,1.7) {$x:0\to1\to2$\\$y:1\to2\to4$，每步乘$2$};\node at (9.95,7.3) {$y=(\frac12)^x$};\begin{scope}[shift={(8.175862068965518,3.3000000000000003)},x=1.6896551724137934cm,y=0.6cm]\draw[->,gray] (-0.4,0)--(2.5,0) node[right]{$x$};\draw[->,gray] (0,-0.5)--(0,5) node[above]{$y$};\draw[orange!80!black,thick,solid,domain=-0.2:2.2,samples=180] plot(\x,{pow(0.5,\x)});\fill[orange!85!black] (0,1.0) circle(2pt);\fill[orange!85!black] (1,0.5) circle(2pt);\fill[orange!85!black] (2,0.25) circle(2pt);\node[below left=5pt] at (0,0){$0$};\node[below=5pt] at (1,0){$1$};\node[below=5pt] at (2,0){$2$};\draw[gray,densely dotted] (1,0)--(1,0.5);\draw[gray,densely dotted] (2,0)--(2,0.25);\draw[->,orange!80!black,thick] (.2,0.8705505632961241)--(.8,0.5743491774985174);\draw[->,orange!80!black,thick] (1.2,0.43527528164806206)--(1.8,0.2871745887492587);\end{scope}\node[align=center,text width=5.8cm] at (9.95,1.7) {$x:0\to1\to2$\\$y:1\to\frac12\to\frac14$，每步减半};
\node[text width=12.6cm,align=center] at (7.1,.65) {两图同尺度；整数点帮助理解作图，任意实数的证明还要看后面的倍数法。};
\end{tikzpicture}\end{center}

\textbf{观察：}横坐标每次增加1，点的高度按倍数变化。\quad\textbf{原因：}a的x+1次方等于a乘a的x次方。

\textbf{考试使用：}先画三个关键点，再连光滑曲线；整数表不代替实数证明。



先选 $a=2$：$2^{-2}=\frac14,\ 2^{-1}=\frac12,\ 2^0=1,\ 2^1=2,\ 2^2=4$。每往右一步就乘 $2$，所以变大。再选 $a=\frac12$：每往右一步就乘 $\frac12$，所以变小。倒数规律还告诉我们 $2^{-x}=(\frac12)^x$，两个表左右倒过来。

\textbf{说明：}这个方法适合\textbf{理解与作图}；只列几个整数，还不足以证明任意实数 $x$ 的结论，下面给出高中数学的规范说明。

\medskip

\end{proofbox}
\clearpage
\begin{proofbox}
\textbf{\textcolor{CProof}{方法二：最快规范证明——比较两个函数值的倍数}}
\begin{center}\begin{tikzpicture}[font=\small,>=stealth,line join=round]
\draw[gray!40,rounded corners=3pt] (0,0) rectangle (14.2,8.8);\node at (7.1,8.200000000000001) {证明的关键：向右走d，函数值恰好乘a的d次方};
\node at (3.1500000000000004,7.3) {$a=2$};\begin{scope}[shift={(1.3391304347826085,3.3142857142857145)},x=2.130434782608696cm,y=0.7857142857142857cm]\draw[->,gray] (-0.3,0)--(2,0) node[right]{$x$};\draw[->,gray] (0,-0.4)--(0,3.8) node[above]{$y$};\draw[blue!75!black,thick,solid,domain=-0.1:1.8,samples=180] plot(\x,{pow(2,\x)});\fill[orange!85!black] (0.5,1.4142135623730951) circle(2pt);\fill[orange!85!black] (1,2) circle(2pt);\node[below=5pt] at (0.5,0){$\frac12$};\node[below=5pt] at (1,0){$1$};\draw[gray,densely dotted] (0.5,0)--(0.5,1.4142135623730951);\draw[gray,densely dotted] (1,0)--(1,2);\end{scope}\node[align=center,text width=5.8cm] at (3.1500000000000004,1.7) {$d=\frac12>0$\\$f(1)=f(\frac12)\,a^{1/2}$};\node at (9.95,7.3) {$a=\frac12$};\begin{scope}[shift={(8.139130434782608,3.3142857142857145)},x=2.130434782608696cm,y=0.7857142857142857cm]\draw[->,gray] (-0.3,0)--(2,0) node[right]{$x$};\draw[->,gray] (0,-0.4)--(0,3.8) node[above]{$y$};\draw[orange!80!black,thick,solid,domain=-0.1:1.8,samples=180] plot(\x,{pow(0.5,\x)});\fill[orange!85!black] (0.5,0.7071067811865476) circle(2pt);\fill[orange!85!black] (1,0.5) circle(2pt);\node[below=5pt] at (0.5,0){$\frac12$};\node[below=5pt] at (1,0){$1$};\draw[gray,densely dotted] (0.5,0)--(0.5,0.7071067811865476);\draw[gray,densely dotted] (1,0)--(1,0.5);\end{scope}\node[align=center,text width=5.8cm] at (9.95,1.7) {$d=\frac12>0$\\$f(1)=f(\frac12)\,a^{1/2}$};
\node[text width=12.6cm,align=center] at (7.1,.65) {两图同尺度；图取d=1/2作示例，正文任取x1<x2，倍数公式适用于任意d>0。};
\end{tikzpicture}\end{center}

\textbf{观察：}同一横向间隔d，对应高度的倍数。\quad\textbf{原因：}同底数幂相除，指数相减。

\textbf{考试使用：}任取两个实数再比较倍数，完成规范单调性证明。


\begin{description}[style=nextline,leftmargin=2.8em,labelsep=0.5em,itemsep=0.55em]
  \item[\textbf{第一步（正值与定点）：}]
  对 $a>0$，正数的实数次幂仍为正，且
  \[
    a^x>0,\qquad f(0)=a^0=1.
  \]
  所以图像在 $x$ 轴上方，并恒过 $(0,1)$。

  \item[\textbf{第二步（先证明 $a>1$ 时递增）：}]
  任取 $x_1<x_2$，则 $d=x_2-x_1>0$。由同底数幂相乘可得
  \[
    \frac{f(x_2)}{f(x_1)}=\frac{a^{x_2}}{a^{x_1}}=a^{x_2-x_1}=a^d.
  \]
  因为 $a>1,\ d>0$，所以 $a^d>1$，又因为 $f(x_1)>0$，得到
  \[
    f(x_2)=f(x_1)a^d>f(x_1).
  \]
  \textbf{为什么 $a^d>1$？}先看 $d=\frac mn>0$：$a^{m/n}=\bigl(\sqrt[n]{a}\bigr)^m>1$。实数指数的运算延续这一大小规律，所以对所有 $d>0$ 都成立。核心是\textbf{从 $x_1$ 到 $x_2$，函数值乘了一个大于 $1$ 的正数}。

  \item[\textbf{第三步（$0<a<1$ 时递减）：}]
  令 $b=\frac1a>1$。当 $d>0$ 时 $b^d>1$，故
  \[
    0<a^d=\frac1{b^d}<1.
  \]
  因而 $f(x_2)=f(x_1)a^d<f(x_1)$。这就说明指数变大，函数值反而变小。

  \item[\textbf{第四步（$x$ 的正负与 $1$ 的关系）：}]
  无论哪种底数，$f(0)=1$。把任意 $x$ 与 $0$ 比较，再利用刚证出的增减性，便得到：
  \[
    \begin{array}{c|cc}
      &x<0&x>0\\\hline
      a>1&0<a^x<1&a^x>1\\
      0<a<1&a^x>1&0<a^x<1
    \end{array}
  \]

\end{description}
\end{proofbox}
\clearpage
\begin{proofbox}

\begin{center}\begin{tikzpicture}[font=\small,>=stealth,line join=round]
\draw[gray!40,rounded corners=3pt] (0,0) rectangle (14.2,8.3);\node at (7.1,7.700000000000001) {越往右越接近0，但每个点仍有正的高度};
\begin{scope}[shift={(1.6660377358490566,2.191176470588235)},x=1.1320754716981132cm,y=3.3823529411764706cm]\draw[->,gray] (-0.5,0)--(4.8,0) node[right]{$x$};\draw[->,gray] (0,-0.16)--(0,1.2) node[above]{$y$};\draw[orange!80!black,thick,dashed,domain=0:4.3,samples=180] plot(\x,{pow(2,-\x)});\fill[orange!85!black] (1,0.5) circle(2.2pt);\fill[orange!85!black] (3,0.125) circle(2.2pt);\fill[orange!85!black] (4,0.0625) circle(2.2pt);\node[below left=5pt] at (0,0){$0$};\node[below=5pt] at (1,0){$1$};\node[below=5pt] at (3,0){$3$};\node[below=5pt] at (4,0){$4$};\node[left=5pt] at (0,1){$1$};\draw[orange!80!black,densely dotted] (1,0)--(1,0.5);\draw[orange!80!black,densely dotted] (3,0)--(3,0.125);\draw[orange!80!black,densely dotted] (4,0)--(4,0.0625);\end{scope}\node[anchor=west,align=left,text width=5.0cm] at (8.4,4.0) {$y=(\frac12)^x$\\[9pt]$x=1:\ y=\frac12$\\[9pt]$x=3:\ y=\frac18$\\[9pt]$x=4:\ y=\frac1{16}>0$\\[9pt]值域$(0,+\infty)$，不含$0$};
\node[text width=12.6cm,align=center] at (7.1,.65) {每多乘一次二分之一，高度更小却仍为正；画在横轴附近，不表示落到了横轴上。};
\end{tikzpicture}\end{center}

\textbf{观察：}小高度仍在横轴上方。\quad\textbf{原因：}每次除以2，正数仍是正数。

\textbf{考试使用：}值域端点0用圆括号，不写最小值为0。

\begin{description}[style=nextline,leftmargin=2.8em,labelsep=0.5em,itemsep=0.55em]
  \item[\textbf{第五步（为什么值域是所有正数）：}]
  当 $a>1$ 时，$a^n$ 随正整数 $n$ 增大而无限增大，而 $a^{-n}=1/a^n$ 无限接近 $0$；$0<a<1$ 时两侧方向相反。指数函数在全体实数上连续，所以会取到 $0$ 与无穷大之间的每个\textbf{正数}，但永远不会等于 $0$。
  \[
  \begin{array}{ll}
    a>1: &\displaystyle\lim_{x\to-\infty}a^x=0,\quad \lim_{x\to+\infty}a^x=+\infty,\\[0.4em]
    0<a<1: &\displaystyle\lim_{x\to+\infty}a^x=0,\quad \lim_{x\to-\infty}a^x=+\infty.
  \end{array}
  \]
  因此值域为 $(0,+\infty)$，$y=0$ 为水平渐近线。
\end{description}

\medskip

\end{proofbox}
\clearpage
\begin{proofbox}
\textbf{\textcolor{CProof}{方法三：用对称点证明图像关于 $y$ 轴对称}}
\begin{center}\begin{tikzpicture}[font=\small,>=stealth,line join=round]
\draw[gray!40,rounded corners=3pt] (0,0) rectangle (14.2,8.3);\node at (7.1,7.700000000000001) {关于y轴对称：横坐标变号，高度和距离都不变};
\begin{scope}[shift={(4.1,2.161111111111111)},x=1.0714285714285714cm,y=0.7301587301587301cm]\draw[->,gray] (-2.8,0)--(2.8,0) node[right]{$x$};\draw[->,gray] (0,-0.7)--(0,5.6) node[above]{$y$};\draw[blue!75!black,thick,solid,domain=-2.3:2.15,samples=180] plot(\x,{pow(2,\x)});\draw[orange!80!black,thick,dashed,domain=-2.3:2.15,samples=180] plot(\x,{pow(2,-\x)});\fill[orange!85!black] (-2,4) circle(2.2pt);\fill[orange!85!black] (2,4) circle(2.2pt);\fill[orange!85!black] (0,1) circle(2.2pt);\node[above=7pt] at (-1.25,4.65){$(-2,4)$};\node[above=7pt] at (1.25,4.65){$(2,4)$};\draw[<->,gray] (-2,4)--(0,4);\draw[<->,gray] (0,4)--(2,4);\node[below=4pt] at (-1,4){$2$};\node[below=4pt] at (1,4){$2$};\node[below=5pt] at (-2,0){$-2$};\node[below left=5pt] at (0,0){$0$};\node[below=5pt] at (2,0){$2$};\end{scope}\node[anchor=west,align=left,text width=5.0cm] at (8.4,4.0) {蓝：$f(x)=2^x$\\[8pt]橙：$g(x)=2^{-x}=f(-x)$\\[8pt]$(2,4)\leftrightarrow(-2,4)$\\[8pt]一般为$(t,f(t))$\\[4pt]$\leftrightarrow(-t,f(t))$};
\node[text width=12.6cm,align=center] at (7.1,.65) {同一水平线上的两点距y轴相等；从一条曲线可直接翻出另一条。};
\end{tikzpicture}\end{center}

\textbf{观察：}同一水平线的两点左右等距。\quad\textbf{原因：}倒数底数等于指数变号。

\textbf{考试使用：}对应点只改横坐标符号，交点为(0,1)。



令 $f(x)=a^x$，$g(x)=(1/a)^x=a^{-x}=f(-x)$。取\textbf{任意}实数 $t$：
\[
  \underbrace{(t,f(t))}_{\text{在 }f\text{ 的图像上}}
  \quad\longleftrightarrow\quad
  \underbrace{(-t,g(-t))}_{\text{在 }g\text{ 的图像上}},
  \qquad g(-t)=f(t).
\]
所以对应的两个点是 $(t,f(t))$ 与 $(-t,f(t))$：\textbf{横坐标互为相反数，纵坐标完全相同}。这正是关于 $y$ 轴对称的条件。

\textbf{易错提醒：}$(t,f(t))$ 关于 $y$ 轴的对称点应写成 $(-t,f(t))$，\textbf{不能}写成 $(-t,f(-t))$；后者的纵坐标一般已经变了。

\medskip
\textbf{\textcolor{CProof}{结论回扣}}

\textbf{三个原因串起整个结论：}正底数的幂始终为正；比较 $a^{x_2}$ 与 $a^{x_1}$ 只需看倍数 $a^{x_2-x_1}$；底数取倒数等同于指数变号。于是定点、增减、值域、渐近线和图像对称性都不必死记。

\end{proofbox}
